\chapter{Analyse et spécification des besoins}
\label{chap:analyse_spec}

\section{Introduction}
La réussite d'un projet informatique, particulièrement dans un contexte industriel complexe et multi-sites comme celui du Groupe COFAT, repose fondamentalement sur une compréhension exhaustive et sans ambiguïté des attentes des futurs utilisateurs. Ce deuxième chapitre est par conséquent consacré à la phase cruciale d'analyse et de spécification des besoins. Dans un premier temps, nous identifierons avec précision les acteurs réels du système, en définissant leurs profils métiers, leurs contraintes opérationnelles, et leurs véritables prérogatives au sein de l'application. Ensuite, nous traduirons ces attentes en exigences concrètes à travers l'énumération détaillée des besoins fonctionnels et non fonctionnels, en veillant à justifier chaque choix par le socle technologique retenu. Afin de visualiser le périmètre d'interaction de chaque acteur, nous présenterons le diagramme de cas d'utilisation global. Puis, nous modéliserons l'architecture statique des données par le biais du diagramme de classes global, reflet fidèle de notre base de données relationnelle. Enfin, nous consoliderons l'ensemble de ces spécifications sous la forme du Backlog Général, qui servira de feuille de route pour la planification des sprints de développement.

\section{Identification des acteurs du système}
L'architecture de l'application \textit{Cofat Capacity Study} repose sur un modèle de gestion des accès basé sur les rôles (RBAC). Cette approche garantit la sécurité des données industrielles tout en offrant une interface utilisateur personnalisée. Les acteurs du système correspondent directement aux entités décisionnelles du groupe COFAT. Trois rôles distincts ont été modélisés :

\subsection{Le Site Manager (Utilisateur principal)}
\textbf{Profil métier :} Le Site Manager est un ingénieur des méthodes ou un responsable de production opérant sur l'un des 10 sites physiques du groupe (par exemple, à Mateur, au Caire, ou à Durango). Il est responsable de la gestion locale de la capacité de son usine.
\newline
\textbf{Contexte et besoins :} Avant ce projet, ce profil passait d'innombrables heures à saisir des données dans des tableurs Excel isolés, rendant difficile la prévision de ses besoins réels face aux nouveaux projets automobiles. Dans la nouvelle plateforme, le Site Manager a besoin d'autonomie. Il doit pouvoir planifier localement ses équipements, gérer ses espaces de production et prévoir ses effectifs. C'est également lui qui crée et soumet les demandes de budget non-industriel pour son site. Il effectue donc un travail de saisie, d'importation de fichiers, d'analyse locale, et de création de demandes budgétaires, mais il n'a pas accès aux données des autres sites du groupe.

\subsection{L'Administrateur (Admin)}
\textbf{Profil métier :} L'Administrateur représente la Direction Centrale ou le Backoffice (généralement basé à Tunis, Kram). Il supervise les opérations globales et gère le système d'information de l'application.
\newline
\textbf{Contexte et besoins :} C'est l'acteur possédant le contrôle total sur le paramétrage du système. Il a la lourde responsabilité de créer les comptes utilisateurs, d'attribuer les rôles, et de gérer le référentiel des sites mondiaux. De plus, il est le principal consommateur des modules de "Consolidation Groupe", ayant accès à la vision agrégée (Dashboard global) des équipements, espaces et ressources humaines mondiaux. \textbf{Note importante concernant le workflow budgétaire :} Conformément à la réalité opérationnelle de l'application, l'Administrateur n'intervient à aucune étape dans le workflow du budget non-industriel. Il n'a aucun pouvoir de validation ou de rejet sur ces demandes.

\subsection{L'Acheteur (Achat)}
\textbf{Profil métier :} Il s'agit d'un membre du département central des Achats, chargé de la négociation avec les fournisseurs externes de COFAT (fabricants de machines, prestataires de services).
\newline
\textbf{Contexte et besoins :} Le rôle de l'Achat a été conçu de manière volontairement ciblée et restreinte dans le système. Sur le module de budget non-industriel, lorsque les Site Managers soumettent une demande (état "Submitted"), l'Achat accède à une liste de ces requêtes. \textbf{Sa seule et unique action est de modifier le coût unitaire (unit price)} d'une ligne de budget en fonction de ses négociations avec les fournisseurs. Cette modification de prix recalcule automatiquement le coût total (total price). L'Achat ne valide pas et ne rejette pas formellement la demande, et n'effectue aucun export de rapport. De même, sur le catalogue \textit{StandardInvestment V2}, l'Achat se limite à la stricte mise à jour des prix/coûts des équipements standards préexistants (par exemple, ajuster le coût d'une machine de coupe Komax suite à l'inflation), sans pouvoir administrer le catalogue (pas de création/suppression d'équipements). Chaque modification de prix déclenche automatiquement une notification au Site Manager concerné.

\section{Spécification des besoins fonctionnels}
Les besoins fonctionnels décrivent avec exactitude le comportement attendu de l'application. Pour chaque fonctionnalité, nous avons pris soin de l'adosser à son implémentation technique dans notre stack.

\begin{itemize}
    \item \textbf{S'authentifier au portail :} Le système doit fournir un écran de connexion. Techniquement, le serveur Node.js génère un token JWT signé (via la bibliothèque jsonwebtoken), stocké dans un cookie HttpOnly côté navigateur, protégeant ainsi l'application contre les attaques XSS. Les mots de passe sont hachés en base de données à l'aide de l'algorithme bcrypt.
    
    \item \textbf{Gérer le référentiel du système (Admin) :} L'application doit permettre à l'Admin d'ajouter de nouveaux utilisateurs, de modifier leurs rôles et de configurer la liste des sites COFAT. L'API Express vérifie le rôle de l'utilisateur à chaque requête via un middleware spécifique avant d'exécuter les opérations de type CRUD sur la base de données SQL Server.
    
    \item \textbf{Importer des données de planification via Excel (Site Manager) :} C'est la fonctionnalité phare pour l'adoption par les usines. L'application doit permettre l'upload d'un fichier .xlsx. Côté frontend (React), l'utilisateur charge le fichier qui est envoyé en format multipart/form-data via Axios. Côté serveur, la bibliothèque \textit{multer} gère le buffer, puis \textit{xlsx} parse les données. Le code valide la structure des colonnes attendues (par exemple, "machineNeed", "availableMachine") et insère massivement les données via la fonction \textit{bulkCreate} de l'ORM Sequelize, en gérant élégamment les lignes vides ou mal formatées.
    
    \item \textbf{Consulter l'analyse de capacité locale (Site Manager) :} L'application doit afficher sous forme de tableaux dynamiques et de graphiques (Charts.js, Recharts) l'état de la charge locale (équipements surchargés, espaces saturés). Le calcul du taux de charge (Load = Besoin / Disponible) est effectué en temps réel.
    
    \item \textbf{Consulter la consolidation mondiale (Admin) :} Le système doit agréger les données locales. Une requête Sequelize optimisée somme les besoins et disponibilités des 10 entités (group by équipement et période) pour offrir une vue consolidée, affichée via des composants React (MUI/Ant Design).
    
    \item \textbf{Gérer le workflow du budget non-industriel :} Le Site Manager doit pouvoir créer une ligne budgétaire. Une fois soumise, l'Achat doit pouvoir la consulter et mettre à jour le coût unitaire (unitPrice). Le système (Node.js) doit automatiquement recalculer le champ \textit{totalPrice} (via un hook Sequelize \textit{beforeUpdate}) lors de cette modification.
    
    \item \textbf{Recevoir des notifications temps réel :} Chaque action sur le budget (création, modification du prix) doit notifier les parties prenantes. L'enregistrement est créé dans la table \textit{BudgetNotification}, et le composant "Cloche" du frontend interroge régulièrement l'API (\textit{polling}) pour afficher le compteur de notifications non lues.
    
    \item \textbf{Générer des exports Excel/PDF :} Le système doit permettre l'export des tableaux consolidés. Node.js génère dynamiquement ces fichiers binaires et les renvoie sous forme de Blob au navigateur.
\end{itemize}

\section{Spécification des besoins non fonctionnels}
Les exigences non fonctionnelles dictent les critères de qualité, de sécurité, et d'infrastructure de l'application. Elles ont largement guidé le choix de l'architecture logicielle :

\begin{itemize}
    \item \textbf{Sécurité des accès et des données :} La protection des stratégies industrielles de COFAT est impérative. Les mots de passe sont salés et hachés avec bcrypt (facteur de coût 10). L'autorisation se fait via un \textit{middleware} Express placé sur chaque route de l'API. Seuls les jetons JWT valides, vérifiés par la clé secrète du serveur, permettent l'accès. Le fait de stocker le jeton en cookie HttpOnly empêche l'exécution de scripts malveillants côté client d'y accéder.
    
    \item \textbf{Performance et fluidité :} Face à l'analyse simultanée de milliers de lignes de données provenant des 10 entités, le système ne doit pas ralentir. L'architecture asynchrone non-bloquante de Node.js est parfaite pour cela. De plus, Microsoft SQL Server assure des temps de réponse d'agrégation de données multi-sites inférieurs à 3 secondes, grâce à une indexation rigoureuse des clés étrangères.
    
    \item \textbf{Extensibilité (Scalabilité) :} Le système devant intégrer progressivement de nouvelles usines ou de nouveaux modules (ex: qualité, maintenance), nous avons opté pour une architecture 3-tiers (Client SPA React, Serveur API REST, BDD relationnelle). Ce découplage strict permet de faire évoluer le backend sans impacter l'interface utilisateur.
    
    \item \textbf{Ergonomie (UI/UX) :} L'application est destinée à des ingénieurs d'usine travaillant souvent sur des tablettes en bord de chaîne, ainsi qu'à des cadres dans des bureaux. Le design doit donc être moderne, clair et responsive. L'utilisation du framework Tailwind CSS garantit l'adaptabilité de l'interface (sans recourir massivement aux media queries CSS), tandis que l'intégration de bibliothèques riches comme Material UI (MUI) et Ant Design permet de manipuler facilement d'immenses tables de données, équipées de filtres et de paginations fluides.
    
    \item \textbf{Portabilité et déploiement continu :} L'application doit pouvoir être installée sur n'importe quel serveur Windows ou Linux du département IT d'ABM Informatique sans conflit de versions (node modules). La solution est donc conçue pour faciliter une future conteneurisation sous Docker et une intégration dans un pipeline CI/CD, garantissant une évolution et des mises en production fluides.
\end{itemize}

\section{Diagramme de cas d'utilisation global}
Afin d'illustrer la répartition des droits et des fonctionnalités entre les trois rôles, nous avons répertorié l'ensemble des cas d'utilisation (Use Cases) dans le tableau \ref{tab:use_cases}, avant de les modéliser via le langage UML.

\begin{table}[h]
\centering
\begin{tabular}{|p{4cm}|p{11cm}|}
\hline
\textbf{Acteur} & \textbf{Cas d'utilisation associés} \\ \hline
\textbf{Tout utilisateur} & - S'authentifier au système \newline - Consulter et gérer ses notifications \\ \hline
\textbf{Administrateur} & - Gérer les utilisateurs (CRUD) \newline - Gérer les sites industriels (CRUD) \newline - Consulter la consolidation Groupe (Équipements, Espaces, RH) \newline - Exporter les données consolidées (Excel, PDF) \\ \hline
\textbf{Site Manager} & - Importer les planifications locales via fichier Excel \newline - Gérer l'analyse locale (Équipements, Espaces, RH) \newline - Créer et soumettre une demande de budget non-industriel \\ \hline
\textbf{Acheteur (Achat)} & - Mettre à jour les coûts/prix des équipements standards (StandardInvestment V2) \newline - Modifier le prix unitaire d'une demande de budget soumise \\ \hline
\end{tabular}
\caption{Liste exhaustive des cas d'utilisation par acteur}
\label{tab:use_cases}
\end{table}

\vspace{0.5cm}
\begin{figure}[H]
    \centering
    % Conserver le nom d'image d'origine pour éviter les erreurs de compilation
    \includegraphics[width=0.9\textwidth]{img/cas_utilisation_global.png} 
    \caption{Diagramme de cas d'utilisation global de Cofat Capacity Study}
    \label{fig:use_case_global}
\end{figure}
\vspace{0.5cm}

\textbf{Analyse du diagramme :}
La figure \ref{fig:use_case_global} met en évidence plusieurs relations fondamentales UML. La relation de dépendance stricte «\textit{include}» est attachée à tous les cas d'utilisation majeurs, pointant vers "S'authentifier". Cela signifie techniquement qu'aucune route de l'API (que ce soit un import Excel ou une mise à jour de prix) ne peut être exécutée sans un jeton de session valide. De même, la relation d'extension «\textit{extend}» est visible : l'action d'importation de fichiers Excel étend la gestion locale des équipements, espaces ou RH, signifiant que la saisie manuelle est la base, mais que l'import en masse vient la compléter en tant qu'option accélérée. Le rôle Achat, quant à lui, est strictement relié à la modification des prix, sans lien direct avec les processus de validation ou d'export.

\section{Diagramme de classes global}
Le modèle de données relationnel, pivot central de l'architecture logicielle, a été conçu en tenant compte de la complexité de l'organisation matricielle du groupe COFAT. Il est directement implémenté dans Microsoft SQL Server via l'ORM Sequelize. Notre schéma s'articule autour de \textbf{11 entités métier} dont les champs reflètent fidèlement le code source réel des modèles Sequelize :

\begin{enumerate}
    \item \textbf{User} : Table \texttt{CofatSearch\_User}. Champs : \texttt{UserId}, \texttt{Username}, \texttt{Email}, \texttt{Password} (haché bcrypt), \texttt{Role} (ENUM : \texttt{admin} | \texttt{user} | \texttt{Achat}), \texttt{CreatedAt}.
    \item \textbf{Site} : Table \texttt{Sites}. Champs : \texttt{id}, \texttt{nom}, \texttt{code}, \texttt{description}, \texttt{pays}, \texttt{actif} (booléen).
    \item \textbf{Equipment} : Table \texttt{Equipment}. Champs : \texttt{id}, \texttt{equipmentId} (unique), \texttt{equipmentCode}, \texttt{nom}, \texttt{referenceEquipment}, \texttt{imagePath}, \texttt{siteId}.
    \item \textbf{EquipmentPlanning} : Table \texttt{EquipmentPlanning}. Champs : \texttt{id}, \texttt{machineNeed} (FLOAT), \texttt{availableMachine} (INT), \texttt{load} (FLOAT), \texttt{toOrder} (INT), \texttt{siteId}, \texttt{equipmentId}, \texttt{year}, \texttt{month}. Une contrainte UNIQUE garantit l'unicité d'une entrée par (siteId, equipmentId, year, month).
    \item \textbf{Space} : Table \texttt{Spaces}. Champs : \texttt{id}, \texttt{type} (zone de production), \texttt{category} (SPACE, Assembly, SUMMARY...), \texttt{year}, \texttt{month}, \texttt{area} (m\textsuperscript{2}), \texttt{rowOrder}, \texttt{siteId}.
    \item \textbf{HR} : Table \texttt{HR}. Champs : \texttt{id}, \texttt{siteId}, \texttt{type} (projet ou zone), \texttt{category} (\texttt{Direct} | \texttt{Assembly Direct} | \texttt{Indirect}), \texttt{year}, \texttt{month} (format \texttt{MO 01}--\texttt{MO 12} ou \texttt{Q 01}--\texttt{Q 04}), \texttt{count} (effectif prévisionnel), \texttt{rowOrder}.
    \item \textbf{NonIndustrialBudget} : Table \texttt{NonIndustrialBudget}. Champs : \texttt{id}, \texttt{department} (IT, HR, QUALITY, BUILDING, LOGISTICS, MAINTENANCE, PRODUCTION), \texttt{area}, \texttt{equipment}, \texttt{qty}, \texttt{currency}, \texttt{unitPrice} (DECIMAL 15,2), \texttt{totalPrice} (DECIMAL 15,2, calculé automatiquement par hook Sequelize), \texttt{createdBy}, \texttt{updatedBy}.
    \item \textbf{BudgetNotification} : Table \texttt{BudgetNotifications}. Champs : \texttt{id}, \texttt{type} (ENUM : NEW\_REQUEST, PRICE\_UPDATED, REQUEST\_VALIDATED, REQUEST\_MODIFIED), \texttt{department}, \texttt{budgetItemId}, \texttt{fromUser}, \texttt{toRole} (ENUM : Achat, User, Admin), \texttt{toUser}, \texttt{message}, \texttt{amount}, \texttt{equipment}, \texttt{isRead}, \texttt{readAt}.
    \item \textbf{Notification} : Alertes système générales pour les opérations sur les modules équipement et espace.
    \item \textbf{StandardEquipment} : Table \texttt{StandardEquipments} (catalogue V1). Champs : \texttt{id}, \texttt{Code\_Eq}, \texttt{Operations}, \texttt{Equipment\_Reference}, \texttt{Supplier\_Technology}, \texttt{Type}, \texttt{Calculation\_Method}, \texttt{Estimated\_Cost\_EUR}, \texttt{Workstation\_Dimensions}, etc.
    \item \textbf{StandardInvestment} : Table \texttt{StandardInvestments} (catalogue V2, version définitive). Champs : \texttt{id}, \texttt{code\_eq}, \texttt{operation}, \texttt{operation\_code}, \texttt{equipment\_reference}, \texttt{supplier\_technology}, \texttt{equipment\_type}, \texttt{calculation\_method}, \texttt{daily\_capacity}, \texttt{lifetime}, \texttt{cost\_euro}, \texttt{cost\_brazil\_real}, \texttt{cost\_mexican\_peso}, \texttt{workstation\_dimensions}, \texttt{reference\_cdc}, \texttt{reference\_pr}, \texttt{QTY}.
\end{enumerate}

\vspace{0.5cm}
\begin{figure}[H]
    \centering
    \includegraphics[width=\textwidth]{img/diagramme_classes_global.png} 
    \caption{Diagramme de classes global de l'application}
    \label{fig:class_diagram_global}
\end{figure}
\vspace{0.5cm}

\textbf{Analyse des relations (Cardinalités et Compositions) :}
Le diagramme illustre six relations structurelles majeures (implémentées en \textit{hasMany} et \textit{belongsTo} dans Sequelize). La classe \textbf{Site} est la clé de voûte de l'architecture multi-tenant :
\begin{itemize}
    \item \textbf{Site $\rightarrow$ EquipmentPlanning (1..N)} : Une usine effectue plusieurs centaines de prévisions sur une année. La suppression d'un site entraînerait logiquement la perte de ses prévisions (composition).
    \item \textbf{Site $\rightarrow$ Spaces (1..N)} et \textbf{Site $\rightarrow$ HR (1..N)} : Les espaces physiques et les effectifs sont strictement inféodés à leur site géographique (Mateur, Durango...).
    \item \textbf{Equipment $\rightarrow$ EquipmentPlanning (1..N)} : Chaque ligne de planning fait obligatoirement référence à une machine spécifique.
    \item \textbf{User $\rightarrow$ Notification (1..N)} : Un employé reçoit plusieurs alertes.
    \item \textbf{NonIndustrialBudget $\rightarrow$ BudgetNotification (1..N)} : L'historique des négociations tarifaires (Achat/Site Manager) est lié à l'article budgétaire concerné.
\end{itemize}

\section{Le Backlog Général (Product Backlog)}
La transformation des spécifications UML en liste de tâches de développement s'est opérée lors des premières réunions avec le Product Owner. Nous avons adopté une priorisation basée sur la "valeur métier immédiate". Autrement dit, les modules permettant de débloquer le processus de planification locale ont reçu la priorité "Haute", la consolidation "Moyenne/Haute", tandis que le workflow budgétaire plus avancé a été planifié en priorité "Moyenne" pour la fin du projet. 

Le tableau \ref{tab:backlog} présente ce Backlog Général, constitué de 14 User Stories (US), qui ont guidé le travail de développement sur la période de mars à septembre 2026.

\begin{table}[h]
\centering
\begin{tabular}{|p{2cm}|p{8cm}|p{2cm}|p{2cm}|}
\hline
\textbf{ID US} & \textbf{Description (En tant que... je veux... afin de...)} & \textbf{Acteur} & \textbf{Priorité} \\ \hline
US-01 & En tant qu'utilisateur, je veux m'authentifier de manière sécurisée afin d'accéder au portail. & Tous & Haute \\ \hline
US-02 & En tant qu'admin, je veux gérer (CRUD) les utilisateurs et leurs rôles. & Admin & Haute \\ \hline
US-03 & En tant qu'admin, je veux paramétrer la liste des sites industriels. & Admin & Haute \\ \hline
US-04 & En tant que Site Manager, je veux importer un fichier Excel de mes équipements locaux. & Site Mgr & Haute \\ \hline
US-05 & En tant que Site Manager, je veux visualiser l'analyse de charge de mes machines (Load \%). & Site Mgr & Haute \\ \hline
US-06 & En tant que Site Manager, je veux importer et analyser mes besoins RH. & Site Mgr & Haute \\ \hline
US-07 & En tant que Site Manager, je veux importer et gérer mes espaces de production. & Site Mgr & Haute \\ \hline
US-08 & En tant qu'admin, je veux visualiser le dashboard consolidé mondial des Équipements. & Admin & Moyenne \\ \hline
US-09 & En tant qu'admin, je veux visualiser la consolidation des Espaces et des RH du groupe. & Admin & Moyenne \\ \hline
US-10 & En tant qu'acheteur, je veux consulter et mettre à jour les coûts unitaires du catalogue StandardInvestment. & Achat & Moyenne \\ \hline
US-11 & En tant que Site Manager, je veux créer et soumettre une demande de budget non-industriel. & Site Mgr & Moyenne \\ \hline
US-12 & En tant qu'acheteur, je veux modifier le prix unitaire d'une demande budgétaire soumise. & Achat & Moyenne \\ \hline
US-13 & En tant qu'utilisateur, je veux recevoir une notification temps réel lors d'une mise à jour sur mon budget. & Tous & Basse \\ \hline
US-14 & En tant que Site Manager ou Acheteur, je veux pouvoir exporter les données de budget en format Excel et PDF. & Tous & Basse \\ \hline
\end{tabular}
\caption{Backlog Général du projet}
\label{tab:backlog}
\end{table}

\textbf{Synthèse de planification :} Ce backlog de 14 User Stories a été réparti logiquement sur les 4 Sprints du projet. Le Sprint 1 a absorbé le socle (US-01 à 03). Le Sprint 2 s'est focalisé sur la planification locale (US-04 à 07). Le Sprint 3 a matérialisé la puissance d'agrégation mondiale et des standards (US-08 à 10). Enfin, le Sprint 4 a bouclé le périmètre avec le workflow de budgétisation collaborative (US-11 à 14).

\section{Conclusion}
Ce deuxième chapitre nous a permis de structurer de manière rigoureuse les fondations du développement logiciel. En clarifiant le rôle exact de chaque acteur — notamment les limitations volontaires apportées au rôle Achat — et en listant exhaustivement les besoins fonctionnels (authentification JWT, import massif Excel, consolidation temps réel) et non fonctionnels (architecture 3-tiers, performances SQL, sécurité), nous avons sécurisé la trajectoire du projet. Les modélisations UML (cas d'utilisation et classes) ont prouvé que la conception de la base de données est parfaitement alignée avec les processus métier de COFAT. Le Backlog Général, fort de ses 14 User Stories, constitue désormais notre contrat de développement.
La prochaine étape, détaillée dans le chapitre suivant, plongera au cœur de la réalisation technique. Nous y argumenterons les choix des technologies et des frameworks, décrirons l'architecture logicielle retenue, avant de présenter les interfaces concrètes du Sprint 1, première brique fonctionnelle du portail web.
